% !PRINTED 432
\weberpage{432}
%<<CONT
durch bilden, dass man die Variable $\omega$ durch einen äquivalenten Werth ersetzt (Bd. I, \S.\,120).

\begin{center}
\rule{0.25\linewidth}{0.4pt}
\end{center}

Wir wollen hier endlich auch noch die Lösung der Ikosaëdergleichung durch hypergeometrische Reihen kurz erwähnen. Nach \gesperrt{Gauss} verstehen wir unter
\[
F(\alpha, \beta, \gamma, \xi)
\]
die unendliche Reihe
\[
1 + \frac{\alpha\beta}{1.\gamma}\xi + \frac{\alpha(\alpha+1)\beta(\beta+1)}{1.2\,\gamma(\gamma+1)}\xi^2 + \frac{\alpha(\alpha+1)(\alpha+2)\beta(\beta+1)(\beta+2)}{1.2.3\,\gamma(\gamma+1)(\gamma+2)}\xi^2 + \cdots
\]
% [sic] The source clearly repeats $\xi^2$ in the third term; the hypergeometric series requires $\xi^3$.
Setzen wir darin
\[
z = 1728\,\xi,
\]
so erhält die Wurzel der Ikosaëdergleichung
\[
H^3 - zf^5 = 0
\]
den Ausdruck
\[
x = \xi^{\frac{1}{3}}\,\frac{F\left(\dfrac{11}{60}, \dfrac{19}{60}, \dfrac{4}{3}, \xi\right)}{F\left(\dfrac{11}{60}, \dfrac{-1}{60}, \dfrac{1}{3}, \xi\right)}.
\]
Diese Reihen convergiren nur, so lange der absolute Werth von $\xi$ ein echter Bruch ist, können aber für andere Werthe von $\xi$ durch ähnliche Reihen ersetzt werden.\footnote{\gesperrt{H. A. Schwarz}, \glqq{}Ueber diejenigen Fälle, in denen die \gesperrt{Gauss}'sche Reihe $F(\alpha,\beta,\gamma,x)$ eine algebraische Function ihres vierten Elementes ist\grqq{}. Crelle's Journal, Bd. 75, 1872 (gesammelte mathematische Werke, Bd. II). \gesperrt{F. Klein}, Vorlesungen über das Ikosaëder, Abschnitt I, Cap. III.}

\begin{center}
\rule{0.25\linewidth}{0.4pt}
\end{center}
